\chapter{From Synchronic Budgets to Diachronic Chains}
\label{ch:synchronic-to-diachronic}

\begin{quotation}
\noindent\textit{Synopsis.} This chapter identifies the conceptual gap that
motivates the monograph's primary contribution. Existing repair-budget
theory describes tradeoffs among distinctions maintained by one observer at
one moment; real information systems instead consist of chains of
transformations performed by many independent agents across time. Local
repair-budget reasoning cannot simply be applied stage by stage, because no
global coordination mechanism exists and each stage acts with incomplete
knowledge of future tasks and other stages' decisions. This motivates the
new formal object introduced in Part~\ref{part:budget-relay}.
\end{quotation}

\section{The Mismatch, Stated Precisely}

\cref{thm:redistribution} governs competition among distinctions
maintained by one observer at one time. The continuation problem of
Chapter~\ref{ch:continuation-problem} concerns a single distinction, or a
bounded cluster of distinctions, whose status must survive transmission
across a sequence of observers or processes over time. These are
different objects, and the difference is not cosmetic.

\section{Why This Is Not a Trivial Relabeling}

Three features of the continuation problem have no counterpart in the
Redistribution Theorem's setting. First, no single global budget
functional constrains an entire chain of independently operating stages;
each stage's budget is its own. Second, each stage's repair decisions are
typically made without visibility into any other stage's decisions, past
or future. Third, the chain is directional: a later stage cannot
un-discard a distinction an earlier stage has already collapsed, whereas
\cref{thm:redistribution}'s variational argument implicitly allows
reallocation in any direction within a single observer's simultaneous
budget.

The temptation to treat these as differences of degree rather than of
kind is worth resisting directly, since a natural first response to the
continuation problem is to try applying \cref{thm:redistribution}
stage-by-stage -- to ask what a ``local Redistribution Theorem'' at each
$B_i$ would say, and to hope that composing $n$ local instances yields
something like a global conservation law across the whole chain. The
remainder of this section shows why that hope fails, and why it fails for
a structural reason rather than merely a technical one.

\begin{proposition}[No Global Conserved Quantity Across a Relay]
\label{prop:no-global-conserved-quantity}
There is, in general, no functional $\Phi_{\mathrm{relay}} =
f(\beta_0, \ldots, \beta_{n-1})$ of a budget relay's stage budgets
satisfying an analogue of \cref{thm:redistribution}'s conservation law
under variation of any single $\beta_i$.
\end{proposition}

\begin{proofsketch}
\cref{thm:redistribution}'s conservation law depends essentially on the
interference functional $I(\cdot,\cdot)$ being defined over pairs of
distinctions competing within a \emph{single} repair-stable domain
$\restdomain$, so that a variation in the repair investment devoted to
one distinction has a well-defined, computable effect on the repair cost
of another via $\interf{\cdot}{\cdot}$ (\cref{def:interference}). Under
\cref{def:budget-relay}, however, stage $i$'s interference function $I_i$
is defined only over distinctions present at stage $i$; no interference
term $I(d, d')$ is defined, or is even well-typed, for $d$ a distinction
competing for $\beta_i$ and $d'$ a distinction competing for $\beta_j$ at
a different stage $j$, since the two distinctions need not both survive
to be simultaneously present at any single stage for such a term to be
evaluated. A variation $\delta\beta_i$ therefore has no defined effect on
$\beta_j$ for $j \neq i$ within this formalism, not merely an
uncoordinated or unobserved one: the mathematical object
$\cref{thm:redistribution}$'s proof varies -- a single $\Phi$ built from
terms sharing one common domain $\restdomain$ -- has no counterpart
spanning stages, since $\restdomain$ itself is stage-local by
\cref{def:substrate-chain}. No conservation law can be stated for a
quantity that has not been defined, and $\Phi_{\mathrm{relay}}$ has not
been defined for exactly this reason.
\end{proofsketch}

\cref{prop:no-global-conserved-quantity} sharpens the informal claim made
in earlier drafts of this argument -- that ``no global coordination
mechanism exists'' -- into a specific mathematical reason rather than a
description of a missing institutional arrangement. It is not merely that
no one has coordinated the stages of a typical relay; it is that the
formal object a coordination claim would need to be about does not exist
under \cref{def:budget-relay}'s stage-local interference structure. This
is also why Chapter~\ref{ch:coordination}'s partial-coordination
mechanisms cannot simply consist of ``computing $\Phi_{\mathrm{relay}}$
and applying \cref{thm:redistribution} to it'': that theorem is
unavailable in this setting not because the chain is uncoordinated but
because the relevant conserved quantity has no definition to conserve.
Coordination mechanisms must therefore do something other than extend
Chapter~\ref{ch:repair-economics}'s apparatus outward; they must
introduce new machinery, which is exactly what
Chapter~\ref{ch:coordination}'s taxonomy does.

\section{Statement of the Extension Task}

The remainder of this monograph is organized around building a formal
object capable of expressing ``locally rational, globally uncoordinated''
loss accumulation -- the budget relay of Chapter~\ref{ch:budget-relay} --
and proving a composition result analogous in spirit to
\cref{thm:redistribution} but diachronic in structure: the Composition
Failure Theorem of Chapter~\ref{ch:composition-failure}. Given
\cref{prop:no-global-conserved-quantity}, that theorem cannot be obtained
by generalizing \cref{thm:redistribution}'s variational argument to a
larger domain; it must be proved directly against
\cref{def:substrate-chain} and \cref{def:budget-discard}'s irreversibility,
as it is in Appendix~\ref{app:proof}.

\begin{principle}[Global Insufficiency, restated]
\label{principle:global-insufficiency-restated}
A substrate chain in which each transformation is locally
budget-rational can nonetheless become globally insufficient for an open
task space, without any single stage being negligent.
\end{principle}

This sentence is, in an important sense, the actual center of the
monograph -- more so than any single theorem statement that follows from
it. It is stated here informally, and precisely, because everything in
Part~\ref{part:budget-relay} exists to make it rigorous, and
\cref{prop:no-global-conserved-quantity} is the first indication that
making it rigorous will require genuinely new machinery rather than a
larger application of what Chapter~\ref{ch:repair-economics} already
provides.
